% P11 paper module: paper-internal rough/logarithmic extension of the odd prime-cell certificate.
% This module repairs the R9 self-containment gap.  It uses only notation and identities
% already established in the manuscript, in particular the prime-cell construction in
% the proof of Theorem~\ref{thm:odd}.

\subsection{Rough prime-cell quadrature and the finite logarithmic gate}
\label{subsec:rough-log-gate}

The smooth proof of Theorem~\ref{thm:odd} used the derivative estimate
\eqref{eq:odd-kernel-bound} in the final prime-cell quadrature.  For the complement
problem below one needs a version that does not assume derivatives of the source vector.
We record that extension here.

Fix $S>0$ and let
\[
 0\ne f\in L^2(-S,S)
\]
be odd and compactly supported in $(-\rho,\rho)$ for some $\rho<S$.  Assume that its
first nonvanishing integral jet
\[
 m=m(f):=\min\{j\ge0:\beta_S^{(j)}(f)\ne0\}
\]
is finite.  Put
\[
 F:=E_Sf,
 \qquad
 \omega_f(h):=\|\tau_hF-F\|_{L^2(\mathbb R)}.
\tag{LG.1}\label{eq:loggate-modulus}
\]
For the future-prime cells of Step~5 in the proof of Theorem~\ref{thm:odd}, write
\[
 \delta_T:=\max_I |I|.
\]
There $q\asymp X=e^{2(T-r)}$ and the fixed short-interval exponent $3/5$ gives
$|I|\asymp X^{-2/5}=e^{-\frac45(T-r)}$.  All future cells used in the certificate
satisfy $T-r\ge T/2-O(1)$; hence
\[
 \boxed{\delta_T\le C e^{-2T/5}.}
\tag{LG.2}\label{eq:loggate-cellwidth}
\]

\begin{lemma}[Rough boundary mean and constant-mode scale]
\label{lem:rough-boundary-mean}
For the fixed compactly supported $L^2$ vector $f$, the boundary expansion of
Theorem~\ref{thm:jet-expansion} extends verbatim in the form
\[
 \langle J_{S,T}f,H_T\mathbf1_T\rangle
 =-\sqrt2\,e^{T/2}T^{-1/2}
 \sum_{j=0}^{M}\frac{c_j}{T^j}\beta_S^{(j)}(f)
 +O_{S,M,\|f\|_2}(e^{T/2}T^{-M-3/2})
\tag{LG.3}\label{eq:rough-boundary-expansion}
\]
for every fixed $M$.  Consequently the growing primitive mean $K_T$ from
\eqref{eq:odd-KT} satisfies
\[
 K_T
 =-\sqrt2\,c_m\beta_S^{(m)}(f)
 \frac{e^{T/2}}{T^{m+1/2}}(1+o(1)),
\tag{LG.4}\label{eq:rough-KT}
\]
and
\[
 M_T:=\frac{|K_T|^2}{2T}
 =c_m^2|\beta_S^{(m)}(f)|^2
 \frac{e^T}{T^{2m+2}}(1+o(1)).
\tag{LG.5}\label{eq:rough-MT}
\]
\end{lemma}

\begin{proof}
The proof of Theorem~\ref{thm:jet-expansion} first obtains a kernel expansion uniform
for $0\le r\le S$.  Since $f\in L^2(-S,S)\subset L^1(-S,S)$ with
$\|f\|_1\le\sqrt{2S}\|f\|_2$, that uniform expansion can be integrated directly
against $f$, giving~\eqref{eq:rough-boundary-expansion}.  The finite primitive block and
all prime powers $k\ge2$ form the bounded remainder of
\eqref{eq:odd-rem-bounded}; its pairing with $\mathbf1_T$ is $O(\sqrt T)$, negligible
against the leading term in~\eqref{eq:rough-boundary-expansion}.  This proves
\eqref{eq:rough-KT}, and~\eqref{eq:rough-MT} is immediate.
\end{proof}

For the growing primitive block keep the notation
\[
 a_p:=\tfrac12\log p,
 \qquad d_p:=T-a_p,
 \qquad c_p:=\sqrt{\log p}\,p^{-3/4},
\]
and define, exactly as in~\eqref{eq:odd-kernel},
\[
 k_T(t):=-2\sum_p c_p f(t-d_p),
 \qquad
 k_T^0(t):=k_T(t)-\frac{K_T}{T}.
\tag{LG.6}\label{eq:rough-kernel}
\]
The sums are restricted to the same growing future-prime block as in the proof of
Theorem~\ref{thm:odd}.

\begin{lemma}[Derivative-free continuous signed cost]
\label{lem:rough-continuous-cost}
For every fixed compactly supported $f$ as above,
\[
 e^{-T}\int_0^T e^{t/2}|k_T(t)|^2\,dt
 \le \frac{C_S}{T}\|f\|_2^2.
\tag{LG.7}\label{eq:rough-weighted-kernel}
\]
After the mean subtraction in~\eqref{eq:rough-kernel}, the continuous signed
certificate of~\eqref{eq:odd-signed-continuous} therefore still satisfies
\[
 \boxed{\|Y_T^{\rm cont,-}\|^2=o(M_T).}
\tag{LG.8}\label{eq:rough-cont-cost}
\]
No derivative of $f$ is used.
\end{lemma}

\begin{proof}
Put $A=T-t$ and $F_1(u):=e^{u/4}f(u)$.  Since $a_p=A+u$ on a translated copy of
$f$, one has
\[
 e^{-A/4}k_T(T-A)
 =-2\sum_p\sqrt{\log p}\,p^{-7/8}F_1(a_p-A).
\tag{LG.9}\label{eq:rough-convolution}
\]
The positive discrete measure in~\eqref{eq:rough-convolution} has total variation
\[
 \sum_{p\le e^{2T}}\sqrt{\log p}\,p^{-7/8}
 \ll \frac{e^{T/4}}{\sqrt T},
\tag{LG.10}\label{eq:rough-seven-eight}
\]
by Chebyshev's bound $\pi(x)\ll x/\log x$ and partial summation (equivalently, a
dyadic decomposition).  Young's inequality for convolution with a finite discrete
measure gives
\[
 \int_0^T e^{-A/2}|k_T(T-A)|^2\,dA
 \le C_S\frac{e^{T/2}}{T}\|f\|_2^2.
\]
Since $e^{-T}e^{t/2}=e^{-T/2}e^{-A/2}$, this is
\eqref{eq:rough-weighted-kernel}.

The continuous certificate estimate in Step~4 of the proof of
Theorem~\ref{thm:odd} depends on $k_T^0$ only through this weighted $L^2$ norm.
Replacing $k_T$ by $k_T^0=k_T-K_T/T$ adds at most
\[
 C\frac{|K_T|^2e^{-T/2}}{T^2}
\]
to the certificate cost.  By~\eqref{eq:rough-MT} this is $o(M_T)$; the contribution
from~\eqref{eq:rough-weighted-kernel} is $O(T^{-1})=o(M_T)$ because
$M_T\to\infty$.  Hence~\eqref{eq:rough-cont-cost} follows.
\end{proof}

\begin{lemma}[Oscillation form of the future-prime quadrature]
\label{lem:rough-prime-quadrature}
Let the mass-normalized coefficients $\lambda_q^{(I)}$ be those of
\eqref{eq:prime-cell-cost}.  If $\Phi:I\to\mathscr H$ is norm-continuous with values in
a Hilbert space, define
\[
 \omega_\Phi(\delta)
 :=\sup_{|r-r'|\le\delta}\|\Phi(r)-\Phi(r')\|.
\]
Then
\[
 \boxed{
 \left\|
 \sum_{q:r_q\in I}\lambda_q^{(I)}\Phi(r_q)
 -\int_I\Phi(r)\,dr
 \right\|
 \le2|I|\,\omega_\Phi(|I|).
 }
\tag{LG.11}\label{eq:rough-cell-quadrature}
\]
For the signed edge field
\[
 \Phi_T^-(r)(x)
 :=2k_T^0(x)\alpha(2r-x)
   -2k_T^0(2r-x)\alpha(x),
\tag{LG.12}\label{eq:rough-edge-field}
\]
one has
\[
 \omega_{\Phi_T^-}(\delta)
 \le C_\alpha\Bigl(
 \delta\|k_T^0\|_2+\omega_{k_T}(2\delta)
 \Bigr),
\tag{LG.13}\label{eq:rough-edge-modulus}
\]
and
\[
 \|k_T\|_2
 \le C\frac{e^{T/2}}{\sqrt T}\|f\|_2,
 \qquad
 \omega_{k_T}(h)
 \le C\frac{e^{T/2}}{\sqrt T}\omega_f(h).
\tag{LG.14}\label{eq:rough-kernel-modulus}
\]
Consequently the source quadrature remainder satisfies
\[
 \boxed{
 \|Z_T^{\rm quad,rough}\|_2
 \le C e^{T/2}\sqrt T
 \Bigl[\delta_T\|f\|_2+\omega_f(2\delta_T)\Bigr]
 +C\sqrt T\,\delta_T|K_T|.
 }
\tag{LG.15}\label{eq:rough-quadrature-bound}
\]
\end{lemma}

\begin{proof}
Because the coefficients are positive and
$\sum_{q:r_q\in I}\lambda_q^{(I)}=|I|$, choose any $r_I\in I$ and subtract
$\Phi(r_I)$ from both the discrete sum and the integral.  The triangle inequality gives
\eqref{eq:rough-cell-quadrature}.

For~\eqref{eq:rough-edge-modulus}, the first term in~\eqref{eq:rough-edge-field}
changes only through the fixed smooth anchor $\alpha$, and is bounded by
$C_\alpha\delta\|k_T^0\|_2$.  The second term changes by a translation of $k_T^0$ by
at most $2\delta$.  The constant part $-K_T/T$ has zero translation modulus, giving
\eqref{eq:rough-edge-modulus}.

From~\eqref{eq:rough-kernel}, Minkowski's inequality gives
\[
 \|k_T\|_2\le2\Bigl(\sum_pc_p\Bigr)\|f\|_2,
 \qquad
 \omega_{k_T}(h)\le2\Bigl(\sum_pc_p\Bigr)\omega_f(h).
\]
Chebyshev plus partial summation yields
\[
 \sum_{p\le e^{2T}}\sqrt{\log p}\,p^{-3/4}
 \ll\frac{e^{T/2}}{\sqrt T},
\]
which proves~\eqref{eq:rough-kernel-modulus}.

Finally sum~\eqref{eq:rough-cell-quadrature} over the future cells.  Their total
$r$-length is $O(T)$, so
\[
 \|Z_T^{\rm quad,rough}\|_2
 \le CT\,\omega_{\Phi_T^-}(\delta_T).
\]
Use~\eqref{eq:rough-edge-modulus},~\eqref{eq:rough-kernel-modulus} and
$\|k_T^0\|_2\le\|k_T\|_2+|K_T|/\sqrt T$ to obtain
\eqref{eq:rough-quadrature-bound}.
\end{proof}

\begin{proposition}[Exact sufficient translation-modulus gate]
\label{prop:rough-modulus-gate}
If
\[
 \boxed{
 \omega_f(2\delta_T)=o(T^{-m-3/2}),
 }
\tag{LG.16}\label{eq:rough-modulus-gate}
\]
then
\[
 \boxed{\|Z_T^{\rm quad,rough}\|_2=o(\sqrt{M_T}).}
\tag{LG.17}\label{eq:rough-quadrature-small}
\]
Under~\eqref{eq:rough-modulus-gate}, the dual-certificate proof of
Theorem~\ref{thm:odd} therefore extends to this fixed compactly supported odd $f$ and
produces
\[
 \boxed{
 \sigma_T(J_{S,T}f)
 =c_m^2|\beta_S^{(m)}(f)|^2
 \frac{e^T}{T^{2m+2}}(1+o(1)).
 }
\tag{LG.18}\label{eq:rough-sharp-asymptotic}
\]
\end{proposition}

\begin{proof}
Divide~\eqref{eq:rough-quadrature-bound} by
$\sqrt{M_T}\asymp_f e^{T/2}T^{-m-1}$ from~\eqref{eq:rough-MT}.  The cell-width term is
bounded by
\[
 C_f T^{m+3/2}\delta_T\longrightarrow0
\]
by~\eqref{eq:loggate-cellwidth}.  The last, constant-anchor term has ratio
$O(T\delta_T)\to0$.  The remaining term has ratio
\[
 C_fT^{m+3/2}\omega_f(2\delta_T)=o(1)
\]
by~\eqref{eq:rough-modulus-gate}.  This proves
\eqref{eq:rough-quadrature-small}.

The mass-normalized certificate-cost identity~\eqref{eq:prime-cell-cost} does not use
derivatives, so Lemma~\ref{lem:rough-continuous-cost} gives the same
$o(M_T)$ discrete certificate cost.  The full-rest lift
\eqref{eq:future-primitive-tail}--\eqref{eq:full-rest-lift-cost} is purely operator
algebra and remains unchanged.  The bounded hub remainder is still
\eqref{eq:odd-rem-bounded}.  Thus the upper squeeze
\eqref{eq:full-hub-cost}--\eqref{eq:odd-upper-sharp} carries over, while the
constant-mode lower bound uses only Lemma~\ref{lem:rough-boundary-mean}.  The two sides
match at~\eqref{eq:rough-MT}, proving~\eqref{eq:rough-sharp-asymptotic}.
\end{proof}

\begin{lemma}[Finite logarithmic regularity implies the required modulus]
\label{lem:log-translation}
For $\alpha>0$ define
\[
 \HG^\alpha(\mathbb R)
 :=\left\{F\in L^2(\mathbb R):
 \int_{\mathbb R}[\log(2+|\xi|)]^{2\alpha}
 |\widehat F(\xi)|^2\,d\xi<\infty\right\}.
\tag{LG.19}\label{eq:log-space}
\]
Then
\[
 \boxed{
 F\in\HG^\alpha
 \Longrightarrow
 \|\tau_hF-F\|_2
 =o\bigl((\log(1/h))^{-\alpha}\bigr)
 \quad(h\downarrow0).
 }
\tag{LG.20}\label{eq:log-translation}
\]
\end{lemma}

\begin{proof}
By Plancherel,
\[
 \|\tau_hF-F\|_2^2
 =\frac1{2\pi}\int|e^{i\xi h}-1|^2|\widehat F(\xi)|^2\,d\xi.
\]
Put $M=h^{-1/2}$.  On $|\xi|\le M$,
$|e^{i\xi h}-1|^2\le h^2\xi^2\le h$, hence the low-frequency contribution is
$O(h\|F\|_2^2)=o((\log(1/h))^{-2\alpha})$.  On $|\xi|>M$ use
$|e^{i\xi h}-1|^2\le4$ and obtain
\[
 \int_{|\xi|>M}|\widehat F(\xi)|^2\,d\xi
 \le [\log(2+M)]^{-2\alpha}
 \int_{|\xi|>M}[\log(2+|\xi|)]^{2\alpha}|\widehat F(\xi)|^2\,d\xi.
\]
The weighted tail tends to zero, while
$\log(2+M)\asymp\log(1/h)$.  Taking square roots proves
\eqref{eq:log-translation}.
\end{proof}

\begin{corollary}[Finite logarithmic threshold]
\label{cor:log-threshold}
If the fixed compactly supported odd vector $f$ has first jet $m$ and
\[
 \boxed{E_Sf\in\HG^{m+3/2}(\mathbb R),}
\tag{LG.21}\label{eq:log-threshold}
\]
then the modulus gate~\eqref{eq:rough-modulus-gate} holds and hence the sharp rough
asymptotic~\eqref{eq:rough-sharp-asymptotic} follows.
\end{corollary}

\begin{proof}
By~\eqref{eq:loggate-cellwidth},
$\log(1/\delta_T)\ge\frac25T-O(1)$.  Lemma~\ref{lem:log-translation} with
$\alpha=m+3/2$ therefore gives
\[
 \omega_f(2\delta_T)=o(T^{-m-3/2}),
\]
which is~\eqref{eq:rough-modulus-gate}.
\end{proof}

\begin{remark}[Scope of the logarithmic gate]
\label{rem:log-gate-firewall}
The condition~\eqref{eq:log-threshold} is a sufficient condition for this robust
prime-cell proof; no necessity is claimed.  The native Gamma form controls only one
power of $\log(2+|\xi|)$ in the quadratic integral, corresponding to the baseline
$\HG^{1/2}$ scale.  Thus a first jet of order $m$ leaves an additional logarithmic
regularity deficit of $m+1$.  Whether the explicit complement vector constructed below
has this additional regularity remains open.  Nothing in this subsection proves or
disproves strong terminal transport.
\end{remark}
